\section{微分係数と導関数}\label{節:微分定義}
前章では,位置から速度を求めるときも,放物線の接線の傾きを求めるときも,平均変化率を作って $h$ を $0$ に近づけた.この操作を,一般の関数について定めよう.ニュートンは,変化する量を「流れる量」,その変化の速度を流率として扱った.時間とともに生成される量を見る考え方である.ここでは変数を時間に限らず,次の極限で微分を定義する.

\begin{definition}{微分係数と導関数}
$f$ が $a$ の近くで定義され,有限の極限
\[
 f'(a)=\lim_{h\to0}\frac{f(a+h)-f(a)}h
\]
が存在するとき,\,$f$ は $a$ で\textbf{微分可能}であるといい,\,$f'(a)$ を $a$ における微分係数という.各点の微分係数を対応させた関数 $f'(x)$ を\textbf{導関数}という.
\end{definition}

微分係数は1つの点での数,導関数は点を動かしたときの関数である.\,$y=f(x)$ のとき $f'(x)$ を $y'$ や $\frac{\dd y}{\dd x}$ とも書く.\,$\frac{\dd y}{\dd x}$ は通常の分数そのものではないが,「何を何で微分しているか」を明示する記号である.

$f(x)=x^2$ なら,
\begin{align*}
 \frac{f(a+h)-f(a)}h
 &=\frac{2ah+h^2}h\\
 &=2a+h\\
 &\longrightarrow2a
\end{align*}
なので,\,$f'(a)=2a$,導関数は $f'(x)=2x$ である.点 $(a,f(a))$ における接線は,傾き $f'(a)$ をもつ直線として
\[
 y-f(a)=f'(a)(x-a)
\]
で表される.前章で求めた接線 $y=2ax-a^2$ は,この式で $f(a)=a^2$,\,$f'(a)=2a$ としたものである.

まず,微分可能なら連続である.実際,\,$h\ne0$ に対して
\[
 f(a+h)-f(a)
 =h\frac{f(a+h)-f(a)}h
\]
であり,右辺は $h\to0$ のとき $0\cdot f'(a)=0$ に近づく.一方,連続だから微分可能とは限らない.\,$f(x)=|x|$ は $0$ で連続だが,差商 $\frac{|h|}{h}$ の右極限は $1$,左極限は $-1$ なので,\,$0$ で微分できない.折れ曲がった点に1つの接線の傾きを与えようとしても,左右の候補が一致しないのだ.

\subsection{微分と一次近似}
微分可能性は,
\[
 f(a+h)=f(a)+f'(a)h+h\varepsilon(h),
 \qquad \varepsilon(h)\to0\quad(h\to0)
\]
と書き直せる.実際,\,$h\ne0$ に対して $\varepsilon(h)=\frac{f(a+h)-f(a)}h-f'(a)$ と置けば,微分係数の定義から $\varepsilon(h)\to0$ となる.\,$\varepsilon(h)$ の値は $h$ によって変わり,正にも負にもなり得る.最後の $h\varepsilon(h)$ は,\,$h$ で割っても $0$ に近づく誤差である.右辺の初めの2項 $f(a)+f'(a)h$ は,点 $(a,f(a))$ における接線の $x=a+h$ での値にあたる.つまり,\,$h$ が小さいとき,関数の値と接線の値の差は $h$ よりもさらに速く $0$ に近づくのである.この形は合成関数の微分を証明するときにも役立つ.
%
